\chapter{Ăn Cũng Xem}
\markboth{Ăn Cũng Xem}{Ăn Cũng Xem}

\begin{center}
	\textit{\color{bookprimary}``Một bữa ăn mà ai cũng nhìn xuống là một bữa ăn no bụng nhưng đói sự hiện diện.''}

	\textit{\color{bookprimary}``A meal in which everyone looks down may fill the stomach while starving presence.''}
\end{center}

\ngancach

\begin{tcolorbox}[colback=bookprimary!6,colframe=bookprimary!35,arc=5pt,title={\bfseries Góc Suy Nghĩ},fonttitle=\bfseries\color{bookprimary}]
Ăn là một hành động rất người.
\textit{(Eating is a deeply human act.)}

Nó không chỉ là nạp thức ăn.
\textit{(It is not only about taking in food.)}

Nó là lúc gia đình thấy nhau, bạn bè nghe nhau, và chính mình cảm nhận được mình đang sống chậm lại một nhịp.
\textit{(It is when families see one another, friends hear one another, and a person feels life slowing down for a moment.)}

Khi điện thoại bò thẳng lên bàn ăn, nó không chỉ lấy mắt.
\textit{(When the phone crawls onto the dining table, it does not only take the eyes.)}

Nó lấy luôn cả cuộc gặp.
\textit{(It takes the meeting itself.)}
\end{tcolorbox}

\section{Bữa Cơm Có Đủ Người Nhưng Thiếu Nhau}
\begin{truyen}{Bữa Cơm Có Đủ Người Nhưng Thiếu Nhau}{A Meal with Everyone Present but Nobody There}
\chuhoa{C}{ó nhà bốn người ngồi ăn.}
\textit{(There is a family of four sitting at dinner.)}

Người cha xem tin tức ngắn.
\textit{(The father watches short news clips.)}

Người mẹ trả lời tin nhắn.
\textit{(The mother replies to messages.)}

Đứa con lớn xem video hài.
\textit{(The older child watches comedy videos.)}

Đứa nhỏ vừa ăn vừa lướt clip hoạt hình.
\textit{(The younger child eats while scrolling cartoon clips.)}

Ai cũng đủ mặt.
\textit{(Everyone is physically there.)}

Không ai thật sự ở cùng ai.
\textit{(No one is truly with anyone else.)}

\adultpair{Nhiều gia đình bây giờ không cãi nhau to. Họ chỉ im lặng cùng màn hình. Kiểu im này nhìn hiền hơn chiến tranh, nhưng nó ăn mòn chậm hơn và sâu hơn. Vì con người vẫn tưởng mình đang ở bên nhau trong khi thật ra đã thôi gặp nhau từ lâu.}{``Many families today do not explode into loud arguments. They simply stay silent with their screens. This kind of silence looks gentler than war, but it erodes more slowly and more deeply. People still think they are together while in truth they stopped meeting one another long ago.''}

Người lớn ngồi đó đôi khi vẫn trách con ít nói.
\textit{(Adults sitting there sometimes still complain that their children do not talk much.)}

Nhưng một đứa trẻ bị nuôi giữa những bữa ăn ai cũng cúi xuống sẽ học được một bài rất rõ.
\textit{(But a child raised among meals where everyone looks down learns a very clear lesson.)}

Bài đó là: màn hình quan trọng hơn con người ở cạnh mình.
\textit{(That lesson is this: the screen matters more than the people beside you.)}
\end{truyen}

\section{Ăn Một Mình Cũng Không Còn Yên}
\begin{truyen}{Ăn Một Mình Cũng Không Còn Yên}{Unable to Eat Alone in Peace}
\chuhoa{C}{ó học sinh ăn sáng một mình nhưng vẫn phải bật video cho có tiếng người.}
\textit{(A student eats breakfast alone yet still needs a video playing for company.)}

Một mình không nhất thiết là cô đơn.
\textit{(Being alone is not necessarily loneliness.)}

Nhưng nếu mọi bữa ăn một mình đều phải có màn hình, con người sẽ quên mất cách ở dịu với chính mình.
\textit{(But if every solitary meal needs a screen, people forget how to stay gently with themselves.)}

\adultpair{Có những bữa ăn đơn sơ nhưng rất lành vì mình thật sự có mặt với nó. Màn hình biến nhiều bữa như vậy thành một ca tiêu thụ nữa, thay vì một khoảnh khắc sống.}{``Some simple meals are deeply healthy because we are truly present with them. The screen turns many such meals into just another act of consumption instead of a moment of living.''}
\end{truyen}

\section{Vừa Nhai Vừa Xem Thành Thói Quen}
\begin{truyen}{Vừa Nhai Vừa Xem Thành Thói Quen}{Chewing and Watching Until It Becomes Habit}
\chuhoa{C}{ó bạn bảo không bật gì lên thì thấy ăn không ngon miệng.}
\textit{(A student says food no longer tastes right unless something is playing on screen.)}

Nghe tưởng vô hại.
\textit{(It sounds harmless.)}

Nhưng đó là lúc một hành động cơ bản của cơ thể đã bị buộc phải đi kèm với kích thích phụ trợ.
\textit{(But that is when a basic bodily act has become dependent on added stimulation.)}

\adultpair{Khi ăn cũng phải xem, đầu óc sẽ dần quên cách nhận ra no, ngon, chậm, và đủ. Nó chỉ học cách nạp cả thức ăn lẫn tiếng ồn cùng lúc.}{``When eating always requires watching, the mind gradually forgets how to recognize fullness, taste, slowness, and enoughness. It learns only to ingest food and noise together.''}
\end{truyen}

\section{Món Ăn Nguội Vì Còn Một Clip}
\begin{truyen}{Món Ăn Nguội Vì Còn Một Clip}{Letting the Food Go Cold for One More Clip}
\chuhoa{C}{ó đĩa cơm bốc khói để đó chỉ vì người ăn còn muốn xem hết một đoạn ngắn.}
\textit{(A hot meal is left cooling because the eater wants to finish one more short clip.)}

Không phải đoạn video đó quý.
\textit{(The clip itself is not precious.)}

Thứ đáng ngại là quyền ưu tiên của đời sống thật đã bị đảo lộn.
\textit{(What is worrying is that the priority order of real life has been reversed.)}

\adultpair{Đời sống hằng ngày được làm nên từ những thứ rất gần: bữa cơm nóng, người ngồi cạnh, một câu hỏi nhỏ. Khi mình để mấy thứ đó thua liên tục trước nội dung rác, mình đang tự làm cùn khả năng quý điều gần gũi.}{``Daily life is made of very near things: a warm meal, the person beside you, one small question. When these keep losing to junk content, we are dulling our own ability to cherish what is close.''}
\end{truyen}

\section{Ngồi Quán Với Bạn Nhưng Mỗi Người Một Góc}
\begin{truyen}{Ngồi Quán Với Bạn Nhưng Mỗi Người Một Góc}{Meeting Friends Yet Retreating into Separate Screens}
\chuhoa{C}{ó nhóm bạn hẹn nhau ăn vặt, gọi món xong rồi cả bàn chìm vào mỗi người một màn hình.}
\textit{(A group of friends meets for snacks, orders, and then sinks into separate screens.)}

Họ không cãi nhau.
\textit{(They are not fighting.)}

Nhưng họ cũng không thật sự gặp nhau.
\textit{(But they are not truly meeting either.)}

\adultpair{Nhiều cuộc gặp bây giờ chỉ còn là việc đặt thân thể gần nhau trong khi sự chú ý thì vắng mặt tập thể. Bạn bè mà chỉ còn chia sẻ địa điểm chứ không chia sẻ hiện diện thì tình bạn cũng mỏng đi rất nhanh.}{``Many meetings today are only bodies placed near each other while attention is collectively absent. If friendship ends up sharing only location and not presence, it thins out very quickly.''}
\end{truyen}

\section{Trẻ Nhỏ Học Cách Ăn Cùng Màn Hình}
\begin{truyen}{Trẻ Nhỏ Học Cách Ăn Cùng Màn Hình}{Younger Children Learning to Eat with Screens}
\chuhoa{C}{ó em bé được dỗ ăn bằng điện thoại từ rất sớm nên lớn lên coi việc nhìn màn hình lúc ăn là bình thường.}
\textit{(A child is fed with a phone from very early on and grows up seeing screen-eating as normal.)}

Nhiều người lớn làm vậy vì mệt hoặc vì tiện.
\textit{(Many adults do this out of exhaustion or convenience.)}

Nhưng cái tiện ngắn của người lớn có thể thành cái lệ thuộc dài của đứa trẻ.
\textit{(But the adult's short convenience can become the child's long dependence.)}

\adultpair{Dạy một đứa trẻ ăn mà không cần màn hình là việc vất vả hơn. Nhưng giáo dục thật thường là như vậy: khó hơn trước mắt để nhẹ hơn về sau.}{``Teaching a child to eat without a screen is harder. But real education is often like that: harder now so that life becomes lighter later.''}
\end{truyen}

\section{Im Lặng Trên Bàn Ăn Không Còn Là Yên Bình}
\begin{truyen}{Im Lặng Trên Bàn Ăn Không Còn Là Yên Bình}{Silence at the Table Is No Longer Peace}
\chuhoa{C}{ó những bữa cơm rất yên, nhưng cái yên đó không phải bình an mà là vì ai cũng bận cúi xuống.}
\textit{(Some meals are very quiet, but that quiet is not peace; it is because everyone is bent downward.)}

Sự im lặng của một gia đình tử tế khác hẳn sự im lặng của bốn người bị tách bởi bốn luồng nội dung khác nhau.
\textit{(The quiet of a healthy family differs greatly from the quiet of four people separated by four streams of content.)}

\adultpair{Đừng nhầm không cãi nhau với đang gần nhau. Có những nhà rất ít ồn nhưng vẫn rất xa. Xa không phải vì ghét, mà vì thói quen không còn trao cho nhau sự chú ý đủ để thành một cuộc gặp.}{``Do not confuse the absence of arguments with closeness. Some homes are quiet yet deeply far apart, not because of hatred but because habit no longer gives each other enough attention to form a meeting.''}
\end{truyen}

\section{Một Bữa Không Máy Là Một Cuộc Gặp}
\begin{truyen}{Một Bữa Không Máy Là Một Cuộc Gặp}{One Screen-Free Meal Is a Real Meeting}
\chuhoa{C}{ó gia đình thử cất hết điện thoại trước bữa tối đầu tiên và thấy khó xử trong mấy phút đầu.}
\textit{(A family tries putting every phone away before dinner and feels awkward in the first few minutes.)}

Rồi sau đó bắt đầu xuất hiện những câu hỏi tưởng nhỏ mà lâu rồi không ai hỏi.
\textit{(Then small questions begin to appear, the sort nobody has asked in a long time.)}

\adultpair{Nhiều nhà không thiếu tình thương. Họ thiếu những khung rất cụ thể để tình thương có chỗ lộ ra. Một bữa ăn không màn hình chính là một khung như vậy.}{``Many families do not lack love. They lack concrete frames in which love can show itself. One screen-free meal is exactly such a frame.''}
\end{truyen}

\begin{baihoc}
Ăn cùng điện thoại quá lâu sẽ làm con người quên mất cảm giác nghe nhau tới hết một câu.
\textit{(Eating with phones for too long makes people forget what it feels like to hear one another through a complete sentence.)}

Một gia đình muốn gần lại đôi khi không cần nói triết lý gì lớn.
\textit{(A family trying to grow closer sometimes does not need any grand philosophy.)}

Họ chỉ cần một cái bàn ăn không có màn hình.
\textit{(They may just need one dining table without screens.)}
\end{baihoc}

\begin{tuhocnhanh}
1. Bữa ăn gần nhất của em có thật sự là một cuộc gặp không? \textit{(Was your most recent meal truly a meeting?)}

2. Trong nhà em, ai là người khó rời điện thoại khỏi bàn ăn nhất? \textit{(In your home, who finds it hardest to keep the phone away from the table?)}

3. Em có dám đề nghị một bữa ăn không màn hình không? \textit{(Would you dare suggest one screen-free meal?)}
\end{tuhocnhanh}

\begin{tongketchuong}
Nhiều gia đình không tan vỡ bằng một trận lớn.
\textit{(Many families do not fall apart through one great explosion.)}

Họ mòn đi bằng những bữa ăn ai cũng có mặt nhưng không ai thật sự ở đó.
\textit{(They wear down through meals in which everyone is present but no one is truly there.)}
\end{tongketchuong}

\ngancach